\section{Cluster decomposition at bounded energy}\label{sec:clusters}

Choose a Kähler form $\omega$ in an integral ample class and fix a positive
integer $\Lambda$. All series in this section are truncated above
$\omega$-degree $\Lambda$. Initially their labels are zero and all integral
homology classes of positive $\omega$-degree. Write $\mathcal J_{\rm isol}$
and $\mathcal J_*$ for the isolated and regular structures of
\cite[Definition~3.13]{DIW}, respectively. At a structure in
$\mathcal J_{\rm isol}$ the connected supports are embedded irreducible
curves, disjoint from each other. At an energy bound, a cluster consists
of a neighborhood $\mathcal U$ of such a support in the Hausdorff space
of compact subsets, a structure $J$ and an embedded core $C$. The
boundary of $\mathcal U$ misses the bounded connected supports, their
classes inside are positive multiples of $[C]$, and $C$ is the unique
curve of that class there \cite[Definition~3.9]{DIW}. A connected
support carries no multiplicity. The cycle insertions below retain it.

\subsection{Excision for disconnected cycles}
Fix $J\in\mathcal J_{\rm isol}$ and let $\mathscr S$ denote the space of nonempty connected supports at $J$
of degree at most $\Lambda$. It is compact by
\cite[Theorem~2.11 and Lemma~3.16]{DIW}. Indeed connectedness is closed
under Hausdorff convergence and every nonempty cycle has integral energy
at least one. Fix a finite partition of $\mathscr S$ into relatively open
and closed subsets $\mathscr S_i$. The union $K_i$ of all curves in
$\mathscr S_i$ is compact by the closed incidence relation in
$\mathscr S_i\times X$. The sets $K_i$ are pairwise disjoint since distinct
irreducible curves are disjoint. Choose open neighborhoods $V_i$ of $K_i$
with disjoint closures. For every sufficiently close tame $J'$ each
bounded connected support is contained in some $V_i$. Otherwise a
sequence of such supports violating the assertion would have a connected
limit at $J$ inside some $K_i$, contrary to Hausdorff convergence. For
a prescribed finite partition by cluster neighborhoods with unoccupied
boundaries, the same argument shows that the partition by the $V_i$
continues to agree with it near $J$. Include the complement of the
selected cluster in the initial partition.

For a union $\mathscr U$ of parts define $P(\mathscr U,J)$ to have
constant term one. Its coefficient at $\beta\ne0$ is the characteristic
function of the cycles in class $\beta$ all of whose connected
components are supported in the corresponding open neighborhoods. Use
the class in $H_2(X;\ZZ)$ under inclusion for cycles in an open target.
These locally constant functions have compact support. To see this over
a short closed path, use compactness at fixed energy and the continuity
of the integral homology label \cite[Theorem~2.11 and Remark~2.13]{DIW}.
Once the $V_i$ are separated, restricting the integration current to
one of them amounts to multiplying by a cutoff which is one on its
closure and zero near the other closures. The nonempty supports converge
there in the Hausdorff topology. Occupation of the several neighborhoods
is locally constant. Thus the insertion for a union of parts extends
over the path, and the interval relation preserves its class in $S$.
These short transports may use smooth segments in charts of tame
structures. Compactness and avoidance of the boundaries hold in a
neighborhood of the given structure, so segments in a smaller convex
coordinate neighborhood join the same endpoints even if the original
isotopy path has only $C^1$ parameter dependence.

At either endpoint, open excision into $\coprod_i V_i$ decomposes each
cycle uniquely as $\sum_i D_i$, including $D_i=0$. In total class
$\beta$ the characteristic function is the sum of external products
over all $\sum_i\beta_i=\beta$. There are finitely many relevant
labels by compactness. Consequently
\begin{equation}\label{eq:excision}
                P(\mathscr U,J)=\prod_{i:\,\mathscr S_i\subset\mathscr U}
                                    P(\mathscr S_i,J).
\end{equation}
Shrinking the $V_i$ about the $K_i$ gives the same insertions by open
excision. The separation in the target is what identifies the right
side with multiplication in $S$.

\subsection{Crossings of the wall}
For a core in class $A$ consider the subgroup of the units of the
truncated series algebra generated by the cluster factors with core
class $dA$, $d\ge2$, at isolated structures. Equivalence modulo this
subgroup is an exact identity after multiplication by finitely many
such factors and their inverses. Every factor has constant term one.

The subgroup is generated by corrections with unobstructed cores.
Indeed, for a cluster around a possibly obstructed core in class $dA$,
keep its outer neighborhood and its complement fixed as above. Choose
a nearby structure in $\mathcal J_*$ in the neighborhood given by
\cite[Propositions~3.14 and~3.21]{DIW}. Transport its insertion along
a short path. At the new endpoint the outer neighborhood partitions
into finitely many clusters with unobstructed cores of classes $kdA$,
$k\ge1$, by Proposition~3.21(2) of that work. Include the complementary
part and apply \eqref{eq:excision} to this partition. This expresses
the entire correction as the product of the resulting regular factors,
even if the former core disappears.

Consider first a transverse fold of class-$A$ cores
\cite[Proposition~3.33]{DIW}. The whole path lies in
$\mathcal J_{\rm isol}$, including the fold parameter, where the
obstructed core is isolated. Spatial separation there extends the
outer cluster insertion across the fold. At the two-core endpoint isolate
the two degree-one branches and partition the remainder. At the other
endpoint there are no class-$A$ cores. Equation~\eqref{eq:excision} at
both ends gives
\[
             P_+P_-\prod_i H_i=\prod_j K_j,
\]
where each $H_i,K_j$ has core class a proper multiple of $A$. Thus the
two distinguished factors are mutual inverses modulo those corrections.
This product accounts also for cycles meeting both core neighborhoods.

Along a regular arc of cores in the isolated locus, the factor is
constant modulo the same subgroup. Use the finite cover by outer
neighborhoods in \cite[Proposition~3.30]{DIW}. On an overlap choose
a common smaller cluster around the core
\cite[Proposition~3.19]{DIW}. Each difference is a union of clusters
in proper multiples, and \eqref{eq:excision} compares the factors.
At a wall along a path fixing the core, use the two-parameter model
$s=x(x-t)$ of \cite[Proposition~3.37]{DIW}. Along $s=0$ the two branches
are $x=0$ and $x=t$. At small $s>0$ the root
$(t+\sqrt{t^2+4s})/2$ joins the distinguished branch at $t=-\delta$
to the other branch at $t=\delta$. For $s<0$ the two positive roots
at the latter end meet at the transverse fold $t=2\sqrt{-s}$.
Use isolated auxiliary paths as in that proposition. The endpoint
insertions persist under the small displacements in $s$ by spatial
separation at the two original endpoints. Regular transport on the
first arc and the fold relation on the second therefore invert the
distinguished factor at the fixed-core crossing.

\subsection{The elementary endpoint}
For an embedded core $C$ with induced complex structure $j$, choose
a spin line $L$ and give $L\oplus L$ the Kähler form
$\Omega=\pi^*(\omega|_C)+i\partial\bar\partial\|w\|^2$ on a small
disk, using equal Hermitian metrics. The symplectic normal bundle
of $C$ has first Chern number $2g-2$. Classification of complex
rank-two bundles over a surface gives a unitary normal isomorphism
with $L\oplus L$. Choose it as the derivative of a tubular embedding
using the symplectic orthogonal to $TC$. The forms agree as bilinear
forms along $C$. Linear interpolation on a smaller disk and the
relative Moser argument, with primitive obtained by fiber contraction,
identify the symplectic forms there. Extend the compatible integrable
structure from a smaller tube by extending its metric and using polar
decomposition with $\omega$. Inside the tube introduce the Lee
deformation with small amplitude $T\chi$, where $\chi$ is one on an
inner disk $U$ and supported in the region of agreement. The resulting
ambient $J_L$ is tame and agrees with the elementary structure on $U$
\cite[Proposition~3.2 and proof of Theorem~7.1]{IP}.

Take $C\subset V\Subset U$ a smaller disk bundle. At $J_L$ no cycle
contained in $\overline V$ meets $\partial V$, since every irreducible
component there lies on $C$. Compactness at bounded energy excludes
such boundary supports for all sufficiently close structures.
Containment in $V$ thus gives a compact open and closed insertion over
any sufficiently short closed path. At $J_L$ it equals the point
insertion of Proposition~\ref{prop:local}. At a nearby generic
structure fixing the unobstructed core, as in
\cite[Proposition~3.54]{DIW}, it is the cluster insertion in a small
Hausdorff ball about $C$. Indeed a sufficiently small ball implies
containment in $V$, and otherwise a sequence of connected supports
in $\overline V$ outside that ball would limit to $C$. This proves
the comparison inside the original compact target.

\begin{proposition}\label{prop:finite}
Through degree $\Lambda$ the image of the compact series in $S$ is a
finite product of elementary factors $\Phi(A_g(Q^\gamma))$ and their
inverses. Each occurrence carries the genus and integral class of its core.
\end{proposition}
\begin{proof}
Start with a generic structure and the finite partition of
\cite[Proposition~3.18]{DIW}. Follow the generic fixed-core path and its
finite subdivision in the proof of \cite[Theorem~3.25]{DIW} to the
elementary endpoint just identified. Regular arcs transport the factor,
and each crossing inverts it modulo the proper-multiple factors.
The sign is $(-1)^w$, where $w$ counts the fixed-core crossings relative to the elementary endpoint. Keep the corrections in the exact identities and replace
each by its regular factors as above. For each recursive correction
the positive integer $\lfloor\Lambda/\langle[\omega],A\rangle\rfloor$
strictly decreases. The branching and the depth are finite. Finally
the insertion for the \emph{entire} $\beta$-locus extends from the
generic structure to the given projective structure over a tame path
by compactness at bounded energy and continuity of the integral class.
The endpoint relation in $S$ completes the proof.
\end{proof}
